\documentclass[11pt]{article}
\usepackage[margin=1in]{geometry}
\usepackage{amsmath,amssymb,amsthm,booktabs,hyperref}
\newtheorem{proposition}{Proposition}
\newtheorem{lemma}{Lemma}
\newcommand{\Prb}{\mathbb P}
\title{Contextual Composition: Constrained Speculative Sampling\\\large A finite integration example and a continuation-information boundary}
\author{}\date{13 November 2026}
\begin{document}\maketitle
\begin{abstract}
A speculative sampler can exactly preserve its supplied target law while
nevertheless implement the wrong interpretation of a constraint. We distinguish
legal-prefix masking from conditioning on eventual valid completion, give a
necessary-and-sufficient finite criterion for their agreement, and instantiate
the accepted/rejected histories as native weighted-witness maps in the existing
context--content construction. A two-token model has legal output in every run
but total variation error $1/14$ relative to the intended conditioned law.
The missing interface datum is continuation success mass, not an additional
rejection-sampling trick. This is a worked integration of established finite
probability and speculative sampling, not a new general contextad theorem.
\end{abstract}

\section{What this example tests}
The applications catalogue lists speculative decoding as an integration problem:
combine draft generation, verification and correction with a declared public
semantics. Merely proving the familiar rejection identity would add little.
Here the target itself changes with the admissible continuation interface.
The issue is whether ``legal next token'' contains enough information to
implement ``sample the original model conditional on eventual success.''
It generally does not.

All sets are finite, all weights rational, and the horizon $T$ fixed. There are
no asymptotic speedup claims, floating-point approximations, irreversible tool
calls, continuous samples, or concurrency semantics. We use the one-token
acceptance/correction mechanism underlying speculative decoding, not a new
implementation of its parallel multi-token scheduler.

The existing independent interaction criterion supplies normalized finite
weighted witnesses, dependent multiplication of weights, and native maps
$f:X\to Y$ satisfying $\sum_{f(x)=y}w_x=v_y$. Its probability/list instance
retains zero-weight witnesses and typed ordered output receipts. Its finite
protocol construction evaluates result-dependent continuations by multiplying
weights over dependent histories. These are the ingredients used below.
The plural-semantics results require that observations be named explicitly;
output-only equality is not automatically trace or cost equality.

\section{A target-faithful sampler with an explicit witness map}
Fix distributions $p,q$ on a common finite alphabet $A$. Here $p$ is the supplied
target and $q$ the proposal. Put
\[
 a_i=\min(p_i,q_i),\qquad e_i=(q_i-p_i)_+,\qquad
 d_i=(p_i-q_i)_+,\qquad \rho=\sum_i e_i=\sum_i d_i.
\]
There is an accepted witness $(\mathrm{acc},i)$ of weight $a_i$ for each $i$.
If $\rho>0$, there is a rejected-and-corrected witness $(\mathrm{rej},i,j)$
of weight $e_i d_j/\rho$. If $\rho=0$, no rejection branch is sampled;
zero-weight points in the declared acceptance carrier are retained.
The emitted-token map sends $(\mathrm{acc},i)$ to $i$ and
$(\mathrm{rej},i,j)$ to $j$.

\begin{lemma}[Weighted native comparison]
The emitted-token map pushes these witnesses to $p$ and is a native weighted
map in the finite probability instance.
\end{lemma}
\begin{proof}
At token $j$, the total mass is
$a_j+\sum_i e_i d_j/\rho=a_j+d_j=p_j$ when $\rho>0$.
When $\rho=0$, $p=q$ and $a_j=p_j$. Summing also proves normalization.
\end{proof}

Operationally, first draw $i\sim q$ and accept with probability
$\min(1,p_i/q_i)$ where $q_i>0$. On rejection draw $j$ from $d/\rho$.
No division at a zero-probability proposal is needed. This is the standard
acceptance/correction identity, credited to the speculative-decoding literature
rather than claimed as a new result \cite{leviathan}.

For each public prefix $h$, repeat this construction with $p_h,q_h$. The
continuation interface is the prefix $hi$ after emitting $i$; draft $i$ is not
committed before verification. The witness map preserves this endpoint and the
single-token content receipt. Its body is the deterministic prefix update.
Dependent accumulation therefore extends the pointwise map to a map on complete
histories. Explicitly, induction over $T$ shows that projection has law
$\prod_{t=0}^{T-1}p_{h_t}(i_{t+1})$.
This is an actual use of the existing native weighted-map and dependent
composition rules, not merely equality at one initial distribution.
If content receipts instead retain drafts, rejection flags, timing or cost,
this projection does not preserve those labels. One must declare a weaker
observation; no equality of those stronger observations is claimed.

\section{Legal masking versus conditioning}
Let $P_h$ be the original model's next-token distribution, defined on every
prefix of length below $T$. Let $L\subseteq A^T$ be the successful complete
outputs. Define the backward success mass
\[
 u(h)=\Prb_P(H_T\in L\mid H_{|h|}=h),\qquad
 u(h)=\sum_i P_h(i)u(hi),\quad u(w)=\mathbf1_L(w)\ (|w|=T).
\]
This is a definition via future model kernels, even at a prefix with zero
probability from the chosen initial state. Assume $u(\epsilon)>0$.
At a prefix with $u(h)>0$, let the admissible alphabet be
$A_h^+=\{i:u(hi)>0\}$. This already uses the stronger notion of positive-mass
extendibility, so failure below is not caused by an avoidable dead end.
Two candidate next-token kernels are
\[
 M_h(i)=\frac{P_h(i)\mathbf1_{A_h^+}(i)}{Z(h)},\quad
 Z(h)=\sum_{j\in A_h^+}P_h(j),\qquad
 C_h(i)=\frac{P_h(i)u(hi)}{u(h)}.
\]
$M$ renormalizes only for local legal continuation; $C$ uses future success
mass. These are separately declared normalized grades. Conditioning is not
being smuggled in as a native weight-preserving map from the original $P$ grade.
In general normalization changes weights, and no such native map exists.

\begin{proposition}[Finite continuation criterion]
The law generated by $C$ is $P(\cdot\mid L)$. The complete-output law of $M$
agrees with it if and only if, at every prefix reachable under the conditioned
law, $u(hi)$ is constant over $i$ with $P_h(i)>0$ and $u(hi)>0$.
\end{proposition}
\begin{proof}
Along a successful word $w$, multiplication of the $C$ kernels telescopes:
\[
 \prod_t C_{h_t}(i_{t+1})=
 \prod_t P_{h_t}(i_{t+1})\frac{u(w)}{u(\epsilon)}
 =P(w)/u(\epsilon).
\]
Unsuccessful words have zero mass. At a reachable prefix, equality of
$M_h(i)$ and $C_h(i)$ for positive legal $i$ is equivalent to
$u(hi)=u(h)/Z(h)$, independent of $i$. This proves sufficiency by induction.
For necessity, equality of complete word laws implies equality of prefix
marginals and their one-step conditional laws at every positive-mass prefix.
All paths admitted by $M$ have positive $P$ mass and end in $L$ (an induction
using $u>0$), so the positive reachable prefix sets agree. Thus the displayed
pointwise equality is necessary. Zero-probability branches impose no condition.
\end{proof}

This is elementary finite conditioning, often described as a Doob transform.
The contribution of this example is its explicit interface interpretation and
reproducible connection to the sampler, not novelty of that probability result.
A speculative sampler supplied $M_h$ returns $M$, not $C$; supplied $C_h$ it
returns $C$. Correct verification cannot repair a misdeclared target.

\section{Exact two-token model}
Take $A=\{a,b\}$, $T=2$, and $L=\{aa,ba,bb\}$. The original kernels are
\[
 P_\epsilon=(1/2,1/2),\qquad P_a=(3/4,1/4),\qquad P_b=(1/4,3/4).
\]
Then $u(a)=3/4$, $u(b)=1$, and $u(\epsilon)=7/8$.
The two legal choices at the root have unequal future success mass.
\begin{center}
\begin{tabular}{lrrr}\toprule
Law & $aa$ & $ba$ & $bb$\\\midrule
Original law restricted to $L$ (unnormalized)&$3/8$&$1/8$&$3/8$\\
Global conditioning $C$&$3/7$&$1/7$&$3/7$\\
Local masking $M$&$1/2$&$1/8$&$3/8$\\\bottomrule
\end{tabular}
\end{center}
Both implementations are always legal. Their total variation distance is
$1/14$, witnessed for instance by the event that the first token is $a$.
Use the deliberately imperfect proposal $q_h=(3/4,1/4)$ at each nonterminal
prefix. Exact accepted/rejected witnesses project to the appropriate row
whether the supplied target is $M$ or $C$. Thus draft mismatch and target
misinterpretation are two distinct issues.

Now alter only $P_a$ to $(1/4,3/4)$. The immediate admissible token sets at all
prefixes are unchanged, as is $P_\epsilon$. Yet the correctly conditioned root
law is $(1/5,4/5)$ instead of $(3/7,4/7)$. Consequently a root-level rule whose
only inputs are $P_\epsilon$ and the immediate legality mask cannot compute the
conditioned root law for both models. This is a precise limited-information
obstruction, not a statement that full-model access cannot solve conditioning.
In a finite model, backward dynamic programming does solve it.

\section{What composes, and what must be retained}
A state exposing only the set of legally enabled next tokens is sufficient to
check step legality but insufficient for this conditioned-law specification.
An enriched interface can expose the residual prefix (or an equivalent finite
state) together with its future kernel and terminal success predicate, from
which $u$ is computed. The local equation $u(h)=\sum_iP_h(i)u(hi)$ is a
certificate of consistency between adjacent interfaces.

For any fixed finite horizon and rational kernels, backward evaluation and
rational arithmetic compute $C$ exactly. Installing the native speculative
witness maps at these grades gives a pointwise replacement stable under all
subsequent finite continuations that depend only on emitted tokens and the
preserved public interface. This follows by multiplying and summing equal
one-step conditional kernels; it does not require the continuation to be
independent of previous emitted tokens. A continuation that observes hidden
proposal history or charges implementation-dependent cost is outside this
interface and need not be preserved.

This example does not prove general higher-order contextual completeness,
a new pseudodistributive law, tractable conditioning for arbitrary language
models, finite-state closure for arbitrary constraints, or any speedup of
single-token rejection. For large autoregressive models, future success mass
may be far more expensive than next-token probabilities. The practical
research problem is to find useful interface certificates or approximations,
with explicit law-error bounds, rather than assume that a legality type
already supplies global conditional probabilities.

\section{Checks, scope and relation to earlier work}
\texttt{check.py} uses only Python's standard library and exact rational
arithmetic. It checks 225 three-token target/proposal pairs (including zero
probabilities); 1,573 positive-success binary horizon-two model/language cases;
the full-law continuation criterion; the two displayed tables; the $1/14$
distance; two target-specific speculative compositions; and the identical-mask
counterexample. Enumeration is regression evidence for the implementation,
not a substitute for the proofs above. Run \texttt{python3 check.py}.
No Lean replay or external scientific benchmark was run.

The local source comparison found speculative decoding only as a prospective
catalogue item, rather than an already developed example, in the inspected
contextual-composition programme sources. This is not an exhaustive novelty
claim. The existing foundational papers supply weighted native maps,
dependent composition and explicitly scoped observations. The new artifact is
a concrete integration example with an information boundary for target
selection. Its conceptual ingredients are established probability and sampling.

\begin{thebibliography}{9}
\bibitem{leviathan} Y. Leviathan, M. Kalman, and Y. Matias,
\emph{Fast Inference from Transformers via Speculative Decoding}, ICML 2023,
PMLR 202:19274--19286, Section 2.3 and Appendix A.1.
\url{https://proceedings.mlr.press/v202/leviathan23a.html}.
\bibitem{O1} \emph{General Independent Interaction Criterion}, programme
manuscript, 1 November 2026, sections ``Discrete semiring resources'' and
``Finite probabilistic protocols.''
\bibitem{O3} \emph{Plural Semantics for Result-Dependent Protocols}, programme
manuscript, section on named output observations and raw histories.
\bibitem{cat} \emph{Potential of applications for the generalized contextads
research programme}, catalogue item 10, now under
\texttt{projects/contextual-composition/applications/potential-applications.md}.
\end{thebibliography}
\end{document}
